\documentclass[11pt,a4paper]{article}
\usepackage[utf8]{inputenc}
\usepackage[T1]{fontenc}
\usepackage[english]{babel}
\usepackage{lmodern}
\usepackage{geometry}
\usepackage{xcolor}
\usepackage{tocloft}
\usepackage{hyperref}
\usepackage{enumitem}
\usepackage{parskip}
\usepackage{booktabs}
\usepackage{longtable}
\usepackage{fancyhdr}
\usepackage{lastpage}
\usepackage{listings}
\usepackage{titlesec}
\geometry{margin=2.15cm}
\setlength{\headheight}{14pt}
\setlength{\emergencystretch}{3em}
\setlist{itemsep=3pt,parsep=1pt,topsep=4pt}
\setlength{\parskip}{4pt plus 1pt}
\definecolor{titleblue}{HTML}{123B5D}
\definecolor{softblue}{HTML}{EDF5FB}
\definecolor{softgray}{HTML}{F5F7FA}
\definecolor{accent}{HTML}{0F6E8C}
\hypersetup{
  colorlinks=true,linkcolor=titleblue,urlcolor=accent,
  pdftitle={Scientific Data Analysis v2.0: English publication edition},
  pdfauthor={interemi},
  pdfsubject={Integrated skill and Scientific Workbench release}
}
\titleformat{\section}{\Large\bfseries\color{titleblue}}{\thesection.}{0.6em}{}
\titleformat{\subsection}{\large\bfseries\color{accent}}{\thesubsection.}{0.6em}{}
\setlength{\cftsecnumwidth}{2.8em}
\tocloftpagestyle{fancy}
\lstdefinestyle{promptstyle}{
  basicstyle=\ttfamily\small,backgroundcolor=\color{softgray},
  frame=single,rulecolor=\color{softblue},breaklines=true,columns=fullflexible
}
\pagestyle{fancy}
\fancyhf{}
\fancyhead[L]{Scientific Data Analysis}
\fancyhead[R]{v2.0 guide: English edition}
\fancyfoot[C]{\thepage\ of \pageref{LastPage}}
\newcommand{\code}[1]{\texttt{\detokenize{#1}}}

\begin{document}
\pagenumbering{gobble}
\begin{titlepage}
\centering
\vspace*{1.0cm}
{\Huge\bfseries\color{titleblue} Integrated Guide\\[0.15cm]
Scientific Data Analysis v2.0\par}
\vspace{0.65cm}
{\Large Skill backend and Scientific Workbench\par}
\vspace{0.8cm}
{\large Local-first product: contracts, guided workflows, run bundles,
artifacts, jobs, previews, support, and dual validation\par}
\vspace{0.8cm}
{\large English publication edition: September 23, 2026\par}
\vspace{0.5cm}
\begin{minipage}{0.93\textwidth}
\small\raggedright
Translation of the preserved private v2.0 guide at commit \code{b9f6176}.
Original TeX SHA-256:
\nolinkurl{c154c0c20a8083a8e1ac8e315d6f2211f496a594a783200b1d78061db7a90701}.
Private source:
\nolinkurl{skills/scientific-data-maintainer/.cache/active_references/guides/guia_scientific_data_analysis_v2_0.tex}.

Recorded results and closure steps are historical, not current verification.
The original release declares installation complete only after Phase 11
post-synchronization validation of the installed skill and app. Historical commands are retained for
traceability: use fresh output directories when reproducing them, without
overwriting earlier evidence. The original remains unchanged privately.
\end{minipage}
\vfill
{\large Prepared with Codex\par}
{\large Copyright 2026 interemi\par}
\end{titlepage}

\clearpage
\pagenumbering{roman}
\fancyfoot[C]{\thepage}
\begin{abstract}
\code{scientific-data-analysis} v2.0 is the first release designed as an
integrated product with Scientific Workbench. The skill remains the
deterministic engine for scripts, JSON contracts, run bundles, artifacts,
warnings, errors, provenance, and original preservation. Scientific Workbench
is the macOS app for planning, guidance, execution, jobs, previews, support
bundles, and confirmation controls.

Scientific execution remains local-first, with read-only originals or copied
inputs. The original guide presents Ollama as the local AI route without a
paid API, with OpenAI, Grok/xAI, and Gemini as optional cloud providers. The
Codex bridge supports extended reasoning, code, and reporting; it does not
replace deterministic capability execution.
\end{abstract}
\tableofcontents
\newpage
\pagenumbering{arabic}
\fancyfoot[C]{\thepage\ of \pageref{LastPage}}

\section{Executive overview}
v2.0 brings together two previously separate pieces:
\begin{itemize}[leftmargin=1.5em]
\item \code{scientific-data-analysis}, with 59 public capabilities and app-readiness contracts.
\item Scientific Workbench, a native macOS app for attaching data, reviewing plans, running workflows, and opening results.
\end{itemize}
The app supports planning and operation. Verifiable operational evidence lives
in jobs, summaries, manifests, artifacts, logs, and run bundles; the release
does not promise automatic scientific truth or mandatory cloud execution.

\section{The skill}
The deterministic backend:
\begin{itemize}[leftmargin=1.5em]
\item Registers public capabilities.
\item Inspects and processes data, documents, FITS, notebooks, tables, and packages.
\item Emits \code{summary.json}, \code{manifest.json}, stdout, stderr, and typed artifacts.
\item Declares \code{PASS}, \code{WARNING}, \code{BLOCKED_CONTROLADO}, \code{FAIL}, or audit status \code{ROTO}.
\item Preserves originals through copies or separate outputs.
\item Blocks cleanly for absent optional backends or unsafe inputs.
\end{itemize}

\section{Scientific Workbench}
The product interface:
\begin{itemize}[leftmargin=1.5em]
\item Loads the installed skill registry.
\item Accepts file or folder attachments.
\item Uses the local planner, skill router, or AI to propose workflows.
\item Executes capabilities in traceable runs.
\item Shows jobs, progress, cancellation, retry, resume, and support bundles.
\item Discovers PNG, PDF, CSV/TSV, Markdown, notebook, JSON, and log previews.
\item Provides guided controls for column selection, optional panels, and visual confirmation.
\end{itemize}

\section{Integration responsibilities}
\begin{longtable}{p{0.25\linewidth}p{0.63\linewidth}}
\toprule
\textbf{Layer} & \textbf{Responsibility} \\
\midrule
Skill & Reproducible scripts, envelopes, run bundles, artifacts, and provenance. \\
Scientific Workbench & UI, planning, command review, processes, artifact indexing, and job recovery. \\
Local/cloud AI & Explanation, planning, and assistance; not the sole scientific source of truth. \\
Codex bridge & Extended reasoning, code, reporting, and filesystem tasks under explicit control. \\
\bottomrule
\end{longtable}

\section{v2.0 JSON contract}
The app accepts legacy, v1.8/v1.9, and v2.0 payloads. Main v2.0 fields:
\begin{longtable}{p{0.30\linewidth}p{0.58\linewidth}}
\toprule
\textbf{Field} & \textbf{Use} \\
\midrule
\code{contract_version} & \code{2.0} for new envelopes; earlier compatibility remains. \\
\code{status} / \code{app_status} & Normalized job/UI state. \\
\code{command} & Redacted command with argv, cwd, and provenance. \\
\code{inputs} / \code{outputs} & Declared inputs and outputs. \\
\code{typed_artifacts} & Primary artifact list for Results. \\
\code{warnings} / \code{errors} & Human-readable and parseable messages. \\
\code{next_actions} & Useful continuation actions. \\
\code{original_modified} & \code{false}, \code{true}, or explained \code{unknown}. \\
\code{app_hints} & Short summary, severity, suggested previews, and tags. \\
\code{job_metadata} & Retry, cancellation, exposure, and estimates. \\
\code{safety} & Read-only/copy policy, secret redaction, and confirmations. \\
\bottomrule
\end{longtable}

\section{Artifact types}
Scientific Workbench recognizes frozen v1.8/v1.9 types and v2.0 additions:
\begin{lstlisting}[style=promptstyle]
summary_json, manifest_json, report_md, preview_png, table_csv,
notebook_ipynb, log_txt, qa_report, handoff_bundle, unknown,
fits_product, fits_visual, metadata_json, preview_pdf,
edited_document, app_preview
\end{lstlisting}
Use \code{unknown} when an artifact cannot be typed accurately. Do not invent
an ambiguous type to make the output appear more complete.

\section{Run bundles}
The preferred format remains:
\begin{lstlisting}[style=promptstyle]
run/
  manifest.json
  summary.json
  stdout.txt
  stderr.txt
  command.txt
  artifacts/
  previews/
  reports/
  tables/
  logs/
  next_steps.md
\end{lstlisting}
The app can discover recursively but should prefer \code{typed_artifacts} and
\code{manifest.json} when present.

\section{Local-first, optional cloud, and Codex bridge}
\subsection{Local-first}
The original guide recommends local Ollama: no paid API key is required,
prompts/context remain on the local machine, and it supports planning,
failure explanation, and result summaries.

\subsection{Optional cloud}
The guide lists OpenAI, Grok/xAI, and Gemini as optional providers for reasoning
or writing. They can incur costs, require user keys, and do not execute
capabilities themselves. Provider errors must be classified without leaking
secrets. These are historical provider descriptions, not a current availability
or pricing verification.

\subsection{Codex bridge}
Codex supports code, reports, refactors, and workspace reasoning. It does not
replace the deterministic backend. Workflows must retain verifiable commands,
outputs, artifacts, and logs.

\section{Safety}
The v2.0 integration requires:
\begin{itemize}[leftmargin=1.5em]
\item Read-only original inputs, with edits confined to copies or separate outputs.
\item No API keys or secrets in logs, summaries, support bundles, or exportable prompts.
\item Job records containing commands, statuses, warnings, errors, and next actions.
\item Preflight and confirmation for optional/GUI backends, with clean blocking when prerequisites are absent.
\item A release block for timeout, raw traceback, secret leakage, false success, or original modification.
\end{itemize}

\section{Workflows recorded in Phase 8}
\begin{longtable}{p{0.27\linewidth}p{0.23\linewidth}p{0.38\linewidth}}
\toprule
\textbf{Family} & \textbf{Status} & \textbf{Interpretation} \\
\midrule
Professional table & \code{PASS} & Parseable run bundle with artifacts. \\
Administrative document & \code{WARNING} & Useful intake with explicit limitations. \\
Inherited notebook & Controlled block & Stops interactive input or side effects before execution. \\
Corrupt container & \code{WARNING} & Recovers available content without false success. \\
Time series & \code{WARNING} & Visible date ambiguity. \\
Sensor/measurement & \code{WARNING} & Visible suspicious ranges. \\
FITS/RGB & \code{WARNING} & Successful RGB case; incomplete WCS warns. \\
Absent optional backend & Controlled block & Clean block with next actions. \\
DOCX roundtrip & \code{PASS} & Edits a confirmed copy, not the original. \\
Keynote/GUI & \code{PASS} & Preflight and confirmed export of a copy. \\
\bottomrule
\end{longtable}
Controlled block denotes the original machine status \code{BLOCKED_CONTROLADO}.

\section{Limits}
\begin{itemize}[leftmargin=1.5em]
\item STILTS/TOPCAT, APT, Keynote, TEAREDUCE, IRAF, and iSTARMOD remain optional, expert, or legacy backends.
\item Li I 6707.8 and SB2 tools are specialized, not general-purpose.
\item GUI workflows need explicit confirmation and copied inputs.
\item Cloud AI does not guarantee scientific accuracy.
\item v2.0 is installed only after Phase 11 synchronizes the skill and validates the installed skill with Scientific Workbench.
\end{itemize}

\section{Historical verification commands}
The original workspace-specific paths are shown through aliases:
\code{$PRIVATE_WORKSPACE} is the original editable workspace;
\code{$DATANALYSIS_PYTHON} is its explicitly selected scientific interpreter.
The external dual-gate script belongs to that workspace and is not supplied by
this source publication. These are historical reproduction steps, not the
current public installation procedure.

Main dual gate:
\begin{lstlisting}[style=promptstyle]
cd "$PRIVATE_WORKSPACE"
./run_v2_0_dual_gate.sh
\end{lstlisting}
Skill check:
\begin{lstlisting}[style=promptstyle]
cd "$PRIVATE_WORKSPACE/skill_work/scientific-data-analysis"
"$DATANALYSIS_PYTHON" scripts/sync_public_surface_docs.py --check
\end{lstlisting}
App check:
\begin{lstlisting}[style=promptstyle]
cd "$PRIVATE_WORKSPACE/ScientificWorkbench"
SKIP_UI_SMOKE=1 RUN_DOCUS_BENCHMARK=0 ./script/run_quality_gate.sh
\end{lstlisting}

\section{Troubleshooting}
\begin{longtable}{p{0.27\linewidth}p{0.61\linewidth}}
\toprule
\textbf{Symptom} & \textbf{Next step} \\
\midrule
Absent optional backend & Accept \code{BLOCKED_CONTROLADO}; inspect preflight and \code{next_actions}. \\
Output conflict & Use a fresh temporary folder or change \code{--outdir}. \\
Notebook with \code{input()} & Execute a copy with declared values or block before execution. \\
Missing preview & Inspect \code{manifest.json}, \code{typed_artifacts}, and \code{previews/}. \\
Cloud provider failure & Fall back to Ollama/local or the skill router; do not leak keys in logs. \\
Ambiguous result & Retain \code{WARNING} and explain the limit instead of silently reporting success. \\
\bottomrule
\end{longtable}

\section{Synchronization}
Final synchronization belongs to Phase 11. Pre-sync gates must pass before
copying the editable skill to the installed location. After copying, repeat
checks on the installed skill and test Scientific Workbench against it.

Do not declare v2.0 installed when post-sync validation fails.

\section{Historical decision}
Phase 8 prepared the editable tree for joint documentation. Phase 9 produced
this guide. Phase 10 validated the dual release candidate. Phase 11 installs
the skill only if post-sync validation of the installed skill and Scientific
Workbench passes without FAIL or ROTO.
\end{document}
